\documentclass[a4paper,11pt]{article}

\usepackage[utf8]{inputenc}
\usepackage[T1]{fontenc}
\usepackage[ngerman]{babel}
\usepackage{geometry}
\usepackage{hyperref}
\usepackage{enumitem}
\usepackage{courier}

\geometry{margin=2.5cm}

\newcommand{\cmd}[1]{\texttt{#1}}

\title{Übungsblatt: Erste Schritte mit der Bash}
\author{}
\date{}

\begin{document}
\maketitle

\section*{Ziel}
In diesem Übungsblatt lernst du grundlegende Konzepte der Bash kennen.
Du arbeitest in der \textbf{git-bash unter Windows}.  
Nicht alle benötigten Befehle werden vorgegeben – nutze Hilfe-Funktionen,
Internetrecherche oder probiere aus.

\section*{Hinweise}
\begin{itemize}
  \item Groß- und Kleinschreibung ist in der Bash wichtig.
  \item Befehle werden mit der \texttt{Enter}-Taste ausgeführt.
  \item Nutze bei Bedarf die Hilfe eines Befehls (z.\,B. mit \texttt{-\,-help}).
  \item Frühere Befehle sind über „Pfeil oben“ erreichbar. 
\end{itemize}
 
\section{Orientierung in der Bash}

\begin{enumerate}[label=\arabic*.]
  \item Starte die \textbf{git-bash}.
  \item Finde heraus, in welchem Verzeichnis du dich gerade befindest.
  \item Liste den Inhalt dieses Verzeichnisses auf.
  \item Wechsle in dein Home-Verzeichnis.
\end{enumerate}

\section{Arbeiten mit Verzeichnissen}

\begin{enumerate}[label=\arabic*.]
  \item Erstelle ein neues Verzeichnis mit dem Namen \texttt{bash-uebung}.
  \item Wechsle in dieses Verzeichnis.
  \item Erstelle darin die folgenden Unterverzeichnisse:
  \begin{itemize}
    \item \texttt{texte}
    \item \texttt{bilder}
    \item \texttt{backup}
  \end{itemize}
  \item Überprüfe, ob die Verzeichnisse korrekt erstellt wurden.
\end{enumerate}

\section{Arbeiten mit Dateien}

\begin{enumerate}[label=\arabic*.]
  \item Wechsle in das Verzeichnis \texttt{texte}.
  \item Erstelle eine leere Datei namens \texttt{notizen.txt}.
  \item Schreibe eine kurze Textzeile in diese Datei (z.\,B. deinen Namen).
  \item Zeige den Inhalt der Datei im Terminal an.
  \item Erstelle eine Kopie der Datei im Verzeichnis \texttt{backup} und nenne sie \texttt{notizen.txt.sik}.
\end{enumerate}

\section{Verschieben und Löschen}

\begin{enumerate}[label=\arabic*.]
  \item Verschiebe die Datei \texttt{notizen.txt} in das Verzeichnis \texttt{backup}.
  \item Benenne die Datei dort in \texttt{notizen\_alt.txt} um.
  \item Lösche die Datei mit der Endung \texttt{sik}.
  \item Lösche das Verzeichnis \texttt{bilder}.
  \item Was passiert, wenn ein Verzeichnis nicht leer ist?
\end{enumerate}

\section{Hilfe und Exploration}

\begin{enumerate}[label=\arabic*.]
  \item Finde heraus, wie du dir eine Kurzbeschreibung eines Befehls anzeigen lassen kannst.
  \item Suche einen Befehl, der dir den Inhalt eines Verzeichnisses detaillierter anzeigt.
  \item Probiere aus, was passiert, wenn du die \texttt{Tab}-Taste beim Tippen eines Befehls drückst.
  \item Beende die git-bash.
\end{enumerate}

\newpage

\section*{Lösungen}

\setcounter{section}{1}

\section{Orientierung in der Bash}
\begin{enumerate}[label=\arabic*.]
  \item Rechte Maustaste im Ordner oder Menü.
  \item \cmd {pwd}
  \item \cmd {ls}
  \item \cmd {cd}  (oder \cmd {cd} $\sim$ ) 
\end{enumerate}


\section{Arbeiten mit Verzeichnissen}
\begin{enumerate}[label=\arabic*.]
  \item \cmd {mkdir bash-uebung}
  \item \cmd {cd bash-uebung}
  \item \cmd {mkdir texte bilder backup}
  \item \cmd {ls}
\end{enumerate}


\section{Arbeiten mit Dateien}
\begin{enumerate}[label=\arabic*.]
  \item \cmd {cd texte}
  \item \cmd {touch notizen.txt}
  \item \cmd {echo "BlaFoo" $>>$ notizen.txt} \\ 
        (oder mit \cmd {nano notizen.txt}, Speichern strg+o, verlassen strg+x ) 
  \item \cmd {cat notizen.txt} \\
  (oder \cmd {less notizen.txt}, verlassen mit „q“)
  \item \cmd {cp notizen.txt ../backup/notizen.txt.sik} \\[-0.2cm]
  
  Geht auch in Einzelschritten: \\
  \cmd {cp notizen.txt ../backup/} \\
  \cmd {cd ../backup} \\ 
  \cmd {mv notizen.txt notizen.txt.sik}
\end{enumerate}


\section{Verschieben und löschen}
\begin{enumerate}[label=\arabic*.]
  \item \cmd {mv notizen.txt ../backup/notizen.txt}
  \item \cmd {mv notizen.txt notizen\_alt.txt}
  \item \cmd {rm notizen.txt.sik}
  \item \cmd {cd ..} \\
  \cmd {rm bilder}
  \item Testen an backup: Fehlermeldung\\
  Lösung: \cmd {rm -r backup} (oder \cmd {rm -rf backup})
\end{enumerate}


\section{Hilfe und Exploration}

\begin{enumerate}[label=\arabic*.]
  \item \cmd {ls -\,-help}
  \item \cmd {ls -lah}
  \item Ergänzung eines Befehls oder Dateinamens, falls möglich.\\
  „TabTab“ ist auch hilfreich!
  \item \cmd {exit} oder \cmd {strg+d}
\end{enumerate}
\end{document}
